% Source: 11.basis.b -- this is chapter 11 of the course, not a second half of
% chapter 10. The part letter `b' is inherited from the source note, hence the
% 11.b.n prefix even though this is the chapter's first section.
\renewcommand{\thetheorem}{11.b.\arabic{theorem}}
\setcounter{theorem}{0}
\setcounter{chapter}{10}

\chapter{Bases of a Vector Space (Part B)}

% Start the section counter at `b' so the heading reads "Section 11.b" and
% agrees with the 11.b.n theorem numbers and with the source note 11.basis.b.
\setcounter{section}{1}
\section{The Dimension Theorem}

\subsection{The Strong Exchange Lemma}

Before we establish the main theorem of this chapter, we formalize a stronger version of the Steinitz Exchange Theorem (Lemma D from the previous section). We proved that an independent list cannot be larger than a spanning list. However, the exact mechanism of the proof actually guarantees something even more powerful: we can precisely substitute the independent vectors into the spanning set.

\begin{lemma}[Lemma D': Strong Steinitz Exchange]
\label{lem:strong_steinitz}
Let $V$ be a vector space over $K$, and suppose that $V = \Sp(v_1, \dots, v_n)$ for some list of vectors $v_1, \dots, v_n \in V$. Let $u_1, \dots, u_m \in V$ be a linearly independent list. 

Then $m \le n$, and moreover, one can delete exactly $m$ vectors from the spanning list $v_1, \dots, v_n$ such that the remaining $n-m$ vectors, together with all the vectors $u_1, \dots, u_m$, still span $V$.
\[
\Sp(u_1, \dots, u_m, v'_1, \dots, v'_{n-m}) = V
\]
(where the $v'$ vectors are the survivors of the original spanning list).
\end{lemma}

\begin{proof}
The fact that $m \le n$ was already proven in Lemma D. To prove the \qt{moreover} part, we simply apply the exact step-by-step substitution procedure described in the proof of Lemma D. 
For steps $j = 1, 2, \dots, m$, we sequentially add $u_j$ to the list, behind the vectors $u_1, \dots, u_{j-1}$ already inserted. At each step the list has $n+1$ vectors and still spans $V$, so it is dependent (Lemma C), and we forcefully eject a redundant $v_k$ vector (Lemma B), which preserves the span and brings the length back to $n$. The ejected vector is never one of the $u$'s, by \cref{lem:drop_index_beyond_head}; this is the general step recorded on \cpageref{rem:steinitz_general_step}. After $m$ steps, we have injected all $m$ independent vectors and deleted exactly $m$ of the original spanning vectors. The resulting mixed list still spans $V$. See also \cpageref{rem:steinitz_stronger_form}, where the same statement is drawn out of the proof of Lemma D directly.
\end{proof}

\subsection{The Existence and Uniqueness of Dimension}

We are now ready to state and prove the foundational theorem of finite-dimensional vector spaces. It guarantees that bases always exist and that they provide an absolute, rigid measure of the \qt{size} of a vector space.

\begin{theorem}[The Basis Theorem]
\label{thm:basis_theorem}
Let $V$ be a finite-dimensional vector space over $K$. Then $V$ has a basis with finitely many elements. 

Moreover, \textbf{every} basis of $V$ is finite, and any two bases of $V$ have exactly the same number of elements.
\end{theorem}

\begin{remark}[Infinite Dimensional Spaces]
This theorem is actually true even if $V$ is not finite-dimensional (this is known as the \qt{Hamel Basis Theorem} and requires Zorn's Lemma). However, in this course, we will only prove the theorem for finite-dimensional spaces. The infinite-dimensional case is rigorously covered in advanced algebra courses (e.g., \textit{Grundstrukturen} in the 2nd semester).
\end{remark}

\subsection{Preparation: Two Lemmas}

For the proof of the theorem we need some preparation. The first lemma is the converse reading of the redundancy lemma of the previous chapter: there, a dependent list was shown to contain a vector lying in the span of its predecessors; here, a list in which \emph{no} vector lies in the span of its predecessors is shown to be independent.

\begin{lemma}[Lemma E: No Predecessor Span Means Independent]
\label{lem:no_predecessor_span}
Let $w_1, \dots, w_\ell \in V$ and assume that
\[
w_j \notin \Sp(w_1, \dots, w_{j-1}) \qquad \text{for all } 1 \le j \le \ell.
\]
Then the list $w_1, \dots, w_\ell$ is linearly independent.
\end{lemma}

\begin{proof}
If $w_1, \dots, w_\ell$ were linearly dependent, then by \cref{lem:targeted_redundancy} there would exist an index $1 \le j \le \ell$ with $w_j \in \Sp(w_1, \dots, w_{j-1})$, contradicting the assumption.
\end{proof}

The second lemma is the one that actually produces a basis. Note that it extracts the basis from \emph{within} the given spanning list, rather than merely asserting that one exists somewhere.

\begin{lemma}[Lemma F: Every Finite Spanning List Contains a Basis]
\label{lem:spanning_list_contains_basis}
Suppose that $\Sp(v_1, \dots, v_n) = V$. Then there exists a subset of $\{v_1, \dots, v_n\}$ which is a basis for $V$.
\end{lemma}

\begin{proof}
The proof is algorithmic and is based on $n$ steps. At each step we consider another vector from the list $v_1, \dots, v_n$ and decide whether or not to drop it from the list. At step number $j$ we consider $v_j$:

\begin{description}
    \item[Step 1.] If $v_1 = 0_V$ we drop $v_1$ from the list. If $v_1 \ne 0_V$ we keep it.
    \item[Step $j$, for $2 \le j \le n$.] If $v_j \in \Sp(v_1, \dots, v_{j-1})$ we drop $v_j$. Otherwise we keep it.
\end{description}

(The two steps apply one and the same test. For $j = 1$ the vectors preceding $v_1$ form the empty list, so $\Sp(v_1, \dots, v_{j-1}) = \Sp(\emptyset) = \{0_V\}$ by \cref{rem:span_of_empty_set}, and \qt{$v_1 \in \Sp(\emptyset)$} says exactly \qt{$v_1 = 0_V$}. This is the degenerate case treated in the solution to \cref{exc:redundancy_j_equals_one}.)

After performing the $n$ steps above we obtain a new, possibly shorter, list
\[
v_{i_1}, v_{i_2}, \dots, v_{i_m}, \qquad \text{where } 1 \le i_1 < i_2 < \dots < i_m \le n.
\]

We claim that $\Sp(v_{i_1}, \dots, v_{i_m}) = V$. Indeed, in any of the steps $1 \le j \le n$ we dropped the vector $v_j$ only if $v_j \in \Sp(v_1, \dots, v_{j-1})$, and by \cref{lem:targeted_redundancy} such a vector may be dropped without changing the span. (That lemma produces one particular index $j$, whereas here the index is whichever one the step is at; but the two-inclusion argument proving its second assertion uses only the property $v_j \in \Sp(v_1, \dots, v_{j-1})$ and never how $j$ was found, so it applies to every index having that property.) Therefore,
\[
\Sp(v_{i_1}, \dots, v_{i_m}) = \Sp(v_1, \dots, v_n) = V.
\]

We also claim that $v_{i_1}, \dots, v_{i_m}$ is linearly independent. Indeed, for every $1 \le k \le m$ the vector $v_{i_k}$ survived its step, which means
\[
v_{i_k} \notin \Sp(v_1, v_2, \dots, v_{i_k - 1}),
\]
and since $v_{i_1}, \dots, v_{i_{k-1}}$ are all taken from among $v_1, \dots, v_{i_k - 1}$, so that $\Sp(v_{i_1}, \dots, v_{i_{k-1}}) \subseteq \Sp(v_1, \dots, v_{i_k - 1})$ by \cref{lem:span_is_smallest} \textbf{(b)}, this gives in particular
\[
v_{i_k} \notin \Sp(v_{i_1}, v_{i_2}, \dots, v_{i_{k-1}}).
\]
By \cref{lem:no_predecessor_span}, the list $v_{i_1}, v_{i_2}, \dots, v_{i_m}$ is linearly independent.

Together with the fact that $\Sp(v_{i_1}, \dots, v_{i_m}) = V$, we get that $v_{i_1}, \dots, v_{i_m}$ is a basis for $V$.
\end{proof}

\subsection{Proof of the Basis Theorem}

Because this theorem is so monumental, we will break its proof down into three distinct, manageable statements.

\subsubsection*{Statement 1: V has a finite basis.}

\begin{proof}
Because we assumed $V$ is finite-dimensional, there must exist some finite list of vectors that spans $V$. Let us call this spanning set $S = \{w_1, \dots, w_k\}$. If $S$ is already linearly independent, it is a basis, and we are done.

If $S$ is linearly dependent, we can extract a basis from it using the following systematic reduction algorithm (often called a \qt{Redundancy Filter}):

\begin{center}
\resizebox{0.9\textwidth}{!}{
    \begin{tikzpicture}[scale=1,
        node distance=1.5cm,
        box/.style={draw, thick, rounded corners, align=center, inner sep=2ex},
        arrow/.style={->, thick, >=stealth}
    ]
        \node[box, fill=blue!10] (start) {\textbf{Spanning List} \\ $w_1, w_2, \dots, w_k$};
        
        \node[box, fill=yellow!10, below=1cm of start] (check) {Read left to right. \\ Is $w_j \in \Sp(w_1, \dots, w_{j-1})$?};
        
        \node[box, fill=red!10, left=1cm of check, yshift=-1.5cm] (drop) {\textbf{Yes:} Drop $w_j$. \\ (Span is preserved by Lemma B)};
        \node[box, fill=green!10, right=1cm of check, yshift=-1.5cm] (keep) {\textbf{No:} Keep $w_j$.};
        
        \node[box, fill=blue!10, below=4cm of start] (final) {\textbf{Final Output} \\ A linearly independent list \\ that still spans $V$ (A Basis!)};
        
        \draw[arrow] (start) -- (check);
        \draw[arrow] (check) -- node[above left] {Redundant} (drop);
        \draw[arrow] (check) -- node[above right] {Independent} (keep);
        \draw[arrow] (drop) |- (final);
        \draw[arrow] (keep) |- (final);
    \end{tikzpicture}
}
\end{center}

Specifically:
\begin{enumerate}
    \item If the zero vector $0_V$ is in the list, drop it.
    \item Read the list from left to right. If a vector $w_j$ is a linear combination of its strictly preceding vectors (i.e., $w_j \in \Sp(w_1, \dots, w_{j-1})$), drop it.
\end{enumerate}
By Lemma B (\cref{lem:targeted_redundancy}), dropping this redundant vector does not change the span of the list. We repeat this until we have checked every vector. The remaining list still spans $V$. Furthermore, by \cref{lem:no_predecessor_span}, it must be linearly independent, because we systematically deleted every vector that could be written as a combination of the vectors preceding it. Thus, the remaining list is a finite basis. This is precisely the algorithm of \cref{lem:spanning_list_contains_basis}, applied to the spanning list $w_1, \dots, w_k$.
\end{proof}

\subsubsection*{Statement 2: Every basis of V is finite.}

\begin{proof}
We prove this by contradiction. Suppose there exists a basis $\mathcal{A}$ of $V$ that is \textit{not} a finite set. 

From Statement 1, we know that $V$ definitely possesses at least one finite basis. Let us call this finite basis $S = \{v_1, \dots, v_n\}$. The size of this finite basis is exactly $n$.

Because $\mathcal{A}$ is infinite, it contains infinitely many linearly independent vectors. Therefore, we can arbitrarily select $n + 100$ distinct vectors from $\mathcal{A}$. Let us call them $u_1, \dots, u_{n+100}$.
Because they belong to a basis, this chosen list of $n+100$ vectors is strictly linearly independent (a subset of an independent set is independent, by \cref{lem:subsets_of_independent_sets}).

We now have two lists:
\begin{itemize}
    \item A spanning list $S$ containing exactly $n$ elements.
    \item An independent list containing $n+100$ elements.
\end{itemize}

By the Steinitz Exchange Theorem (Lemma D, \cref{thm:steinitz}), the number of elements in any independent list must be less than or equal to the number of elements in any spanning list. Therefore, we must have:
\[
n + 100 \le n
\]
This implies $100 \le 0$, which is a blatant mathematical contradiction! Thus, our initial assumption must be wrong. Every basis of $V$ must be finite.
\end{proof}

\subsubsection*{Statement 3: Any two bases have the same cardinality.}

\begin{proof}
Let $\mathcal{A}$ and $\mathcal{B}$ be two arbitrary bases of $V$. By Statement 2, we know that both $\mathcal{A}$ and $\mathcal{B}$ must be finite sets. Let $|\mathcal{A}|$ denote the number of elements in $\mathcal{A}$, and $|\mathcal{B}|$ denote the number of elements in $\mathcal{B}$.

Because they are bases, both sets are simultaneously independent and spanning. This allows us to apply Lemma D twice in alternating directions:

\begin{center}
\begin{tikzpicture}[
    node distance=1.5cm,
    box/.style={draw, thick, rounded corners, align=center, inner sep=2ex},
    arrow/.style={->, thick, >=stealth}
]
    \node[box, fill=gray!10] (A) {\textbf{Basis $\mathcal{A}$}};
    \node[box, fill=gray!10, right=4cm of A] (B) {\textbf{Basis $\mathcal{B}$}};
    
    \draw[arrow] (A.north east) to[bend left=20] node[midway, above] {Independent in Spanning} node[midway, below] {$|\mathcal{A}| \le |\mathcal{B}|$} (B.north west);
    
    \draw[arrow] (B.south west) to[bend left=20] node[midway, below] {Independent in Spanning} node[midway, above] {$|\mathcal{B}| \le |\mathcal{A}|$} (A.south east);
\end{tikzpicture}
\end{center}

\begin{enumerate}
    \item Because $\mathcal{A}$ is linearly independent and $\mathcal{B}$ spans $V$, Lemma D (\cref{thm:steinitz}) dictates that $|\mathcal{A}| \le |\mathcal{B}|$.
    \item Reversing their roles, because $\mathcal{B}$ is linearly independent and $\mathcal{A}$ spans $V$, Lemma D dictates that $|\mathcal{B}| \le |\mathcal{A}|$.
\end{enumerate}

Since $|\mathcal{A}| \le |\mathcal{B}|$ and $|\mathcal{B}| \le |\mathcal{A}|$, it must be that $|\mathcal{A}| = |\mathcal{B}|$. All bases contain exactly the same number of elements.
\end{proof}

\subsection{The Concept of Dimension}

Because the number of elements in a basis is a rigid, absolute invariant of the vector space, we can finally define the space's dimension.

\begin{definition}[Dimension]
\label{def:dimension}
Let $V$ be a finite-dimensional vector space over $K$. We define the \textbf{dimension} of $V$ to be the number of elements in any basis of $V$: if $\mathcal{B}$ is a basis of $V$ and $n$ is the number of its elements, we set
\[
\dim V := n \in \mathbb{Z}_{\ge 0}.
\]

If $V$ is not finite-dimensional, we write $\dim V = \infty$.
\end{definition}

\begin{importantremark}
\label{rem:dimension_well_defined}
$\dim V$ is \textbf{well defined}: the definition does not depend on the specific choice of a basis for $V$. This is exactly what Statement 3 of \cref{thm:basis_theorem} buys us. Without it, the phrase \qt{the number of elements in any basis} would not name a single number, and the definition above would be empty.
\end{importantremark}

\begin{remark}[The Zero Space]
\label{rem:dimension_of_zero_space}
Consider the zero vector space $V = \{0_V\}$. It has exactly two subsets, $\emptyset$ and $\{0_V\}$, and both of them span it: $\Sp(\{0_V\}) = \{0_V\} = V$, and $\Sp(\emptyset) = \{0_V\} = V$ as well, by \cref{rem:span_of_empty_set}. Of the two, however, $\{0_V\}$ is linearly dependent, since a list containing the zero vector never is (\cref{lem:independence_basics} \textbf{(d)}). The only linearly independent subset that spans $V$ is therefore the empty set $\emptyset$, which is independent by \cref{rem:empty_list_independent}.
Therefore, the basis of the zero space is $\emptyset$. Because the empty set contains $0$ elements, we have:
\[
\dim \{0_V\} = 0
\]
\end{remark}

\subsection{Examples of Dimensions}

Now that we have formally defined dimension, we can readily determine the dimensions of our familiar vector spaces by simply counting the elements in their standard bases.

\begin{example}[The Coordinate Space]
The standard basis for $K^n$ is $\{e_1, \dots, e_n\}$. Because this set contains exactly $n$ elements, we have:
\[
\dim K^n = n
\]
\end{example}

\begin{example}[Bounded Polynomials]
The standard basis for the space of polynomials of degree $\le d$ is $\{1, x, x^2, \dots, x^d\}$. Be careful when counting here! Because we start at $x^0 = 1$, this set contains $d+1$ elements. Thus:
\[
\dim K[x]_d = d + 1
\]
\end{example}

\begin{example}[Matrices]
The standard basis for the matrix space $\M_{m \times n}(K)$ consists of the individual single-entry matrices $E_{i,j}$. Since there are $m$ rows and $n$ columns, there are exactly $m \cdot n$ such matrices. Thus:
\[
\dim \M_{m \times n}(K) = m \cdot n
\]
\end{example}

\begin{example}[Function Spaces]
Let $S$ be a \emph{finite} set. Recall from \cref{ex:function_space} that the space of all functions from $S$ to $K$, denoted as $K^S := \{f \mid f: S \to K\}$, forms a vector space over $K$.

The dimension of this space matches the cardinality of the underlying set $S$:
\[
\dim K^S = |S|
\]

\begin{exercise}[Basis for the Space of Functions on a Finite Set]
\label{exc:basis_of_function_space}
If $S$ is a finite set, say $S = \{s_1, \dots, s_n\}$, can you construct an explicit basis for $K^S$? \\
\textit{(Hint: Try to define a family of functions $f_i$ where each function outputs $1$ for exactly one specific element $s_i$, and $0$ for all other elements).}
\end{exercise}

\begin{remark}
\label{rem:infinite_function_space}
When the domain $S$ is an infinite set, the space $K^S$ is infinite-dimensional. In this case, we unfortunately do not have a \qt{nice} or easily describable basis for it, and the formula above no longer applies: by \cref{def:dimension} we simply write $\dim K^S = \infty$. It is worth seeing why the obvious candidate fails. The functions $f_s$ of the exercise above are still linearly independent for infinite $S$, but they no longer span $K^S$, because a linear combination is a \emph{finite} sum and so can only produce functions that vanish outside a finite set. \Cref{exc:k_infinity_independent} is exactly this phenomenon for $S = \mathbb{Z}_{\ge 0}$.
\end{remark}
\end{example}

% Creator: Opus 5
\subsection{Solutions to the Exercises}
\label{sec:solutions_bases_part_b}

\begin{exercisesolution}[Solution to \cref{exc:basis_of_function_space}]
Let $S = \{s_1, \dots, s_n\}$ be finite, and follow the hint: for each $1 \le i \le n$ define $f_i \in K^S$ by
\[
f_i(s_j) :=
\begin{cases}
1 & \text{if } j = i, \\
0 & \text{if } j \ne i.
\end{cases}
\]
We claim that $\mathcal{B} := \{f_1, \dots, f_n\}$ is a basis of $K^S$.

\textbf{Spanning.} Let $f \in K^S$ be any function and put $a_i := f(s_i) \in K$. For every $1 \le j \le n$,
\[
\left( \sum_{i=1}^{n} a_i f_i \right)(s_j) = \sum_{i=1}^{n} a_i f_i(s_j) = a_j = f(s_j),
\]
because every term with $i \ne j$ vanishes. The two functions therefore agree at every element of $S$, and so they are equal by \cref{def:equality_of_functions}. Hence $f = \sum_{i=1}^{n} a_i f_i \in \Sp(\mathcal{B})$.

\textbf{Independence.} Suppose $a_1 f_1 + \dots + a_n f_n$ is the zero function. Evaluating at $s_j$ and using the same computation gives $a_j = 0$, and this for every $1 \le j \le n$. So the only linear combination producing the zero function is the trivial one.

Note also that the $f_i$ are pairwise distinct, since $f_i(s_i) = 1 \ne 0 = f_j(s_i)$ whenever $i \ne j$, so $\mathcal{B}$ really has $n$ elements and not fewer. By \cref{def:dimension},
\[
\dim K^S = n = |S|,
\]
which is the formula stated in the example. The construction is the same one that gives the standard basis $e_1, \dots, e_n$ of $K^n$: taking $S = \{1, 2, \dots, n\}$ turns $f_i$ into $e_i$.
\end{exercisesolution}

\begin{ainote}
The notes prepare the Basis Theorem with two lemmas the transcript had dropped,
\cref{lem:no_predecessor_span} and \cref{lem:spanning_list_contains_basis}, both
restored above with their proofs. The first of them closes a real gap: the filter
that extracts a basis only ever tests a vector against the ones \emph{preceding}
it, whereas \cref{prop:dependence_is_redundancy}, which the transcript cited
here, is about a vector lying in the span of \emph{all} the others.

Two cautions about what the exchange procedure of \cref{lem:strong_steinitz}
does and does not maintain. Ejecting a redundant $v_k$ preserves the span and
returns the length to $n$; it says nothing about independence, which may
perfectly well still fail, and does whenever the original spanning list was
itself redundant. What is maintained is that the $u$'s standing at the head are
independent, and that is exactly what lets
\cref{lem:drop_index_beyond_head} keep the ejected vector among the $v$'s.

And the example on function spaces holds only for finite $S$. For infinite $S$
the space $K^S$ is not finite-di\-men\-sional, so \cref{def:dimension} assigns it
$\infty$ and the equation $\dim K^S = |S|$ does not typecheck: the functions
$f_s$ stay independent but stop spanning, a linear combination being a finite
sum. The remark following the example carries that case.
\end{ainote}
